% 《20 瓦的计算机》——通俗版（“随身读”）。
% 每一章 paperback/NN.tex 对应完整版中编号相同的一章 ../chapters/NN_*.tex；
% 在网站上，读者可以在“简版”和“完整版”之间切换每一章。
\documentclass[11pt,openany]{book}
\usepackage[paperwidth=13cm,paperheight=20cm,top=1.8cm,bottom=1.8cm,inner=1.6cm,outer=1.4cm]{geometry}
% Chinese edition: build with XeLaTeX (see the Makefile)
\usepackage[UTF8,heading=true]{ctex}
\usepackage{amsmath,amssymb}
\usepackage{xcolor}
\usepackage[unicode,colorlinks=true,linkcolor=blue!60!black]{hyperref}
\usepackage{microtype}
\input{paperback/pencil} % minimal colour-pencil illustrations, one per chapter
\setlength{\parskip}{0.3em}
\setlength{\emergencystretch}{2em} % narrow paperback page: allow a little extra stretch
\renewcommand{\chaptermark}[1]{\markboth{\small #1}{}}
% neuron type code: first four letters of the primary mediator (as in the full edition)
\newcommand{\nt}[1]{\textsf{\textbf{#1}}}
% pointer to the full edition at the end of each chapter
\newcommand{\fullversion}[1]{\par\bigskip\noindent\small\emph{公式、模型、习题和参考文献见完整版第~#1~章。}\normalsize}

\title{\Huge 20 瓦的计算机\\[14pt] \Large 大脑如何计算——简明易懂版}
\hypersetup{pdftitle={20 瓦的计算机. 大脑如何计算——简明易懂版},pdfauthor={Konstantin Kladko}}
\author{\Large 康斯坦丁·克拉季科}
\date{}

\begin{document}
\frontmatter
\input{paperback/cover} % cover: a collage of the chapter drawings
% copyright page (same licence as the full edition)
\thispagestyle{empty}
\vspace*{\fill}
\noindent{\footnotesize \copyright~2026 Konstantin Kladko（康斯坦丁·克拉季科）。\\[4pt]
本书的正文和插图采用知识共享“署名—非商业性使用—相同方式共享 4.0 国际”许可协议（CC BY-NC-SA 4.0）发布：\url{https://creativecommons.org/licenses/by-nc-sa/4.0/deed.zh-hans}。在注明作者并以相同许可协议发布衍生作品的前提下，可以为非商业目的自由复制、用于教学、翻译和改编本书。\\[4pt]
完整版、源文件和模型：\\\url{https://github.com/kladkogex/20-watt-computer}。}
\clearpage
\tableofcontents
\input{paperback/00_preface}
\mainmatter
\input{paperback/01}
\input{paperback/02}
\input{paperback/03}
\input{paperback/04}
\input{paperback/05}
\input{paperback/06}
\input{paperback/07}
\input{paperback/08}
\input{paperback/09}
\input{paperback/10}
\input{paperback/11}
\input{paperback/12}
\input{paperback/13}
\input{paperback/14}
\input{paperback/15}
\input{paperback/16}
\input{paperback/17}
\input{paperback/18}
\input{paperback/19}
\input{paperback/20}
\input{paperback/21}
\input{paperback/22}
\end{document}
